% Chunk: §3.3 Symmetries of the S-Matrix (part a: Lorentz + internal symmetries)
% Source: Weinberg QFT Vol I, printed pp. 116-124 (PDF pp. 139-147)
% Status: transcribed and source-checked

equation:
\begin{align}
T_{\alpha\beta}^{\mathrm{in/out}\,*}-T_{\beta\alpha}^{\mathrm{in/out}}
={}&-\int d\gamma\,T_{\gamma\beta}^{\mathrm{in/out}\,*}T_{\gamma\alpha}^{\mathrm{in/out}}
\nonumber\\
&\quad\cdot\left([E_\alpha-E_\gamma\pm i\epsilon]^{-1}
-[E_\beta-E_\gamma\mp i\epsilon]^{-1}\right).
\tag{3.2.9}\label{eq:3.2.9}
\end{align}
To prove the orthonormality of the `in' and `out' states, divide Eq.~\eqref{eq:3.2.9}
by $E_\alpha-E_\beta\pm2i\epsilon$. This gives
\begin{align*}
\left(\frac{T_{\alpha\beta}^{\mathrm{in/out}}}{E_\beta-E_\alpha\mp2i\epsilon}\right)^*
+\frac{T_{\beta\alpha}^{\mathrm{in/out}}}{E_\alpha-E_\beta\pm2i\epsilon}
={}&-\int d\gamma\,
\left(\frac{T_{\gamma\beta}^{\mathrm{in/out}}}{E_\beta-E_\alpha\mp i\epsilon}\right)^*
\frac{T_{\gamma\alpha}^{\mathrm{in/out}}}{E_\alpha-E_\beta\pm i\epsilon}.
\end{align*}
The $2\epsilon$s in the denominators on the left-hand side can be replaced with
$\epsilon$s, since the only important thing is that these are positive infinitesimals.
We see then that $\delta(\beta-\alpha)+T_{\beta\alpha}^{\mathrm{in/out}}/(E_\beta-E_\alpha
\pm i\epsilon)$ is unitary. With \eqref{eq:3.1.17}, this is just the statement that the
in and out states form two orthonormal sets of state-vectors. The unitarity of the
$S$-matrix can be proved in a similar fashion by multiplying \eqref{eq:3.2.9} with
$\delta(E_\beta-E_\alpha)$ instead of $(E_\alpha-E_\beta\pm2i\epsilon)^{-1}$.

\subsection{\texorpdfstring{Symmetries of the $S$-Matrix}{Symmetries of the S-Matrix}}

In this section we will consider both what is meant by the invariance of the
$S$-matrix under various symmetries, and what are the conditions on the Hamiltonian
that will ensure such invariance properties.

\subsubsection*{Lorentz Invariance}

For any proper orthochronous Lorentz transformation $x\to\Lambda x+a$, we may
define a unitary operator $U(\Lambda,a)$ by specifying that it acts as in
Eq.~\eqref{eq:3.1.1} on either the `in' or the `out' states. When we say that a theory is
Lorentz-invariant, we mean that the same operator $U(\Lambda,a)$ acts as in
\eqref{eq:3.1.1} on both `in' and `out' states. Since the operator $U(\Lambda,a)$ is
unitary, we may write
\[
S_{\beta\alpha}=\OutBra{\beta}\InKet{\alpha}
=\OutBra{U(\Lambda,a)\beta}\InKet{U(\Lambda,a)\alpha},
\]
so using \eqref{eq:3.1.1}, we obtain the Lorentz invariance (actually, covariance)
property of the $S$-matrix: for arbitrary Lorentz transformations
$\Lambda^\mu{}_{\nu}$ and translations $a^\mu$,
\begin{align}
&S_{p'_1\sigma'_1n'_1;p'_2\sigma'_2n'_2;\ldots,\,
p_1\sigma_1n_1;p_2\sigma_2n_2;\ldots}
\nonumber\\
&=\exp\!\left(ia_\mu(p_1^\mu+p_2^\mu+\cdots-p_1^{\prime\mu}-p_2^{\prime\mu}-\cdots)\right)
\nonumber\\
&\quad\cdot\sqrt{\frac{(\Lambda p_1)^0(\Lambda p_2)^0\cdots
(\Lambda p'_1)^0(\Lambda p'_2)^0\cdots}
{p_1^0p_2^0\cdots p_1^{\prime0}p_2^{\prime0}\cdots}}
\nonumber\\
&\quad\cdot\sum_{\bar\sigma_1\bar\sigma_2\ldots}
D^{(j_1)}_{\bar\sigma_1\sigma_1}\!\left(W(\Lambda,p_1)\right)
D^{(j_2)}_{\bar\sigma_2\sigma_2}\!\left(W(\Lambda,p_2)\right)\cdots
\nonumber\\
&\quad\cdot\sum_{\bar\sigma'_1\bar\sigma'_2\ldots}
D^{(j'_1)*}_{\bar\sigma'_1\sigma'_1}\!\left(W(\Lambda,p'_1)\right)
D^{(j'_2)*}_{\bar\sigma'_2\sigma'_2}\!\left(W(\Lambda,p'_2)\right)\cdots
\nonumber\\
&\quad\cdot S_{\Lambda p'_1\bar\sigma'_1n'_1;\Lambda p'_2\bar\sigma'_2n'_2;\ldots,\,
\Lambda p_1\bar\sigma_1n_1;\Lambda p_2\bar\sigma_2n_2;\ldots}.
\tag{3.3.1}\label{eq:3.3.1}
\end{align}
(Primes are used to distinguish final from initial particles; bars are used to
distinguish summation variables.) In particular, since the left-hand side is
independent of $a^\mu$, so must be the right-hand side, and so the $S$-matrix
vanishes unless the four-momentum is conserved. We can therefore write the part
of the $S$-matrix that represents actual interactions among the particles in the
form:
\begin{equation}
S_{\beta\alpha}-\delta(\beta-\alpha)
=-2\pi iM_{\beta\alpha}\delta^4(p_\beta-p_\alpha).
\tag{3.3.2}\label{eq:3.3.2}
\end{equation}
(However, as we will see in the next chapter, the amplitude $M_{\beta\alpha}$
itself contains terms that involve further delta function factors.)

Eq.~\eqref{eq:3.3.1} should be regarded as a definition of what we mean by the Lorentz
invariance of the $S$-matrix, rather than a theorem, because it is only for certain
special choices of Hamiltonian that there exists a unitary operator that acts as
in \eqref{eq:3.1.1} on both `in' and `out' states. We need to formulate conditions on the
Hamiltonian that would ensure the Lorentz invariance of the $S$-matrix. For this
purpose, it will be convenient to work with the operator $S$ defined by
Eq.~\eqref{eq:3.2.4}:
\[
S_{\beta\alpha}=\bra{\beta}S\ket{\alpha}_0.
\]
As we have defined the free-particle states in Chapter 2, they furnish a
representation of the Poincar\'e group, so we can always define a
unitary operator $U_0(\Lambda,a)$ that induces the transformation \eqref{eq:3.1.1} on
these states:
\begin{align*}
&U_0(\Lambda,a)\ket{p_1,\sigma_1,n_1;p_2,\sigma_2,n_2;\ldots}_0
\nonumber\\
&=\exp\!\left(-ia_\mu(p_1^\mu+p_2^\mu+\cdots)\right)
\sqrt{\frac{(\Lambda p_1)^0(\Lambda p_2)^0\cdots}{p_1^0p_2^0\cdots}}
\nonumber\\
&\quad\cdot\sum_{\bar\sigma_1\bar\sigma_2\ldots}
D^{(j_1)}_{\bar\sigma_1\sigma_1}\!\left(W(\Lambda,p_1)\right)
D^{(j_2)}_{\bar\sigma_2\sigma_2}\!\left(W(\Lambda,p_2)\right)\cdots
\ket{\Lambda p_1,\bar\sigma_1,n_1;\Lambda p_2,\bar\sigma_2,n_2;\ldots}_0.
\end{align*}
Eq.~\eqref{eq:3.3.1} will thus hold if this unitary operator commutes with the
$S$-operator:
\[
U_0(\Lambda,a)^{-1}S\,U_0(\Lambda,a)=S.
\]
This condition can also be expressed in terms of infinitesimal Lorentz
transformations. Just as in Section~2.4, there will exist a set of Hermitian
operators, a momentum $\mathbf{P}_0$, an angular momentum $\mathbf{J}_0$, and a
boost generator $\mathbf{K}_0$, that together with $H_0$ generate the infinitesimal
version of inhomogeneous Lorentz transformations when acting on free-particle
states. Eq.~\eqref{eq:3.3.1} is equivalent to the statement that the $S$-matrix is
unaffected by such transformations, or in other words, that the $S$-operator
commutes with these generators:
\begin{equation}
[H_0,S]=[\mathbf{P}_0,S]=[\mathbf{J}_0,S]=[\mathbf{K}_0,S]=0.
\tag{3.3.3}\label{eq:3.3.3}
\end{equation}
Because the operators $H_0$, $\mathbf{P}_0$, $\mathbf{J}_0$, and $\mathbf{K}_0$
generate infinitesimal inhomogeneous Lorentz transformations on the free states,
they automatically satisfy the commutation relations \eqref{eq:2.4.18}--\eqref{eq:2.4.24}:
\begin{align}
[J_0^i,J_0^j]&=i\epsilon_{ijk}J_0^k, &&\tag{3.3.4}\label{eq:3.3.4}\\
[J_0^i,K_0^j]&=i\epsilon_{ijk}K_0^k, &&\tag{3.3.5}\label{eq:3.3.5}\\
[K_0^i,K_0^j]&=-i\epsilon_{ijk}J_0^k, &&\tag{3.3.6}\label{eq:3.3.6}\\
[J_0^i,P_0^j]&=i\epsilon_{ijk}P_0^k, &&\tag{3.3.7}\label{eq:3.3.7}\\
[K_0^i,P_0^j]&=iH_0\delta_{ij}, &&\tag{3.3.8}\label{eq:3.3.8}\\
[J_0^i,H_0]&=[P_0^i,H_0]=[P_0^i,P_0^j]=0, &&\tag{3.3.9}\label{eq:3.3.9}\\
[K_0^i,H_0]&=iP_0^i. &&\tag{3.3.10}\label{eq:3.3.10}
\end{align}
where $i,j,k$, etc. run over the values 1, 2, and 3, and $\epsilon_{ijk}$ is the
totally antisymmetric quantity with $\epsilon_{123}=+1$.

In the same way, we may define a set of `exact generators,' operators
$\mathbf{P}$, $\mathbf{J}$, $\mathbf{K}$ that together with the full Hamiltonian
$H$ generate the transformations \eqref{eq:3.1.1} on, say, the `in' states. (As already
mentioned, what is not obvious is that the same operators generate the same
transformations on the `out' states.) The group structure tells us that these
exact generators satisfy the same commutation relations:
\begin{align}
[J^i,J^j]&=i\epsilon_{ijk}J^k, &&\tag{3.3.11}\label{eq:3.3.11}\\
[J^i,K^j]&=i\epsilon_{ijk}K^k, &&\tag{3.3.12}\label{eq:3.3.12}\\
[K^i,K^j]&=-i\epsilon_{ijk}J^k, &&\tag{3.3.13}\label{eq:3.3.13}\\
[J^i,P^j]&=i\epsilon_{ijk}P^k, &&\tag{3.3.14}\label{eq:3.3.14}\\
[K^i,P^j]&=iH\delta_{ij}, &&\tag{3.3.15}\label{eq:3.3.15}\\
[J^i,H]&=[P^i,H]=[P^i,P^j]=0, &&\tag{3.3.16}\label{eq:3.3.16}\\
[K^i,H]&=iP^i. &&\tag{3.3.17}\label{eq:3.3.17}
\end{align}

In virtually all known field theories, the effect of interactions is to add an
interaction term $H_{\mathrm{int}}$ to the Hamiltonian, while leaving the momentum
and angular momentum unchanged:
\begin{equation}
H=H_0+H_{\mathrm{int}},\qquad \mathbf{P}=\mathbf{P}_0,\qquad
\mathbf{J}=\mathbf{J}_0.
\tag{3.3.18}\label{eq:3.3.18}
\end{equation}
(The only known exceptions are theories with topologically twisted fields, such as
those with magnetic monopoles, where the angular momentum of states depends on the
interactions.) Eq.~\eqref{eq:3.3.18} implies that the commutation relations \eqref{eq:3.3.11},
\eqref{eq:3.3.14}, and \eqref{eq:3.3.16} are satisfied provided that the interaction commutes with
the free-particle momentum and angular-momentum operators
\begin{equation}
[H_{\mathrm{int}},\mathbf{P}_0]=[H_{\mathrm{int}},\mathbf{J}_0]=0.
\tag{3.3.19}\label{eq:3.3.19}
\end{equation}
It is easy to see from the Lippmann--Schwinger equation \eqref{eq:3.1.16} or equivalently
from \eqref{eq:3.1.13} that the operators that generate translations and rotations when
acting on the `in' (and `out') states are indeed simply $\mathbf{P}_0$ and
$\mathbf{J}_0$. Also we easily see that $\mathbf{P}_0$ and $\mathbf{J}_0$ commute
with the operator $U(t,t_0)$ defined by Eq.~\eqref{eq:3.2.6}, and hence with the
$S$-operator $U(\infty,-\infty)$. Further, we already know that the $S$-operator
commutes with $H_0$, because there are energy-conservation delta functions in both
terms in \eqref{eq:3.2.7}. This leaves just the boost generator $\mathbf{K}_0$ which we need
to show commutes with the $S$-operator.

On the other hand, it is \textit{not} possible to set the boost generator
$\mathbf{K}$ equal to its free-particle counterpart $\mathbf{K}_0$, because then
Eqs.~\eqref{eq:3.3.15} and \eqref{eq:3.3.8} would give $H=H_0$, which is certainly not true in the
presence of interactions. Thus when we add an interaction $H_{\mathrm{int}}$ to
$H_0$, we must also add a correction $\mathbf{W}$ to the boost generator:
\begin{equation}
\mathbf{K}=\mathbf{K}_0+\mathbf{W}.
\tag{3.3.20}\label{eq:3.3.20}
\end{equation}
Of the remaining commutation relations, let us concentrate on Eq.~\eqref{eq:3.3.17}, which
may now be put in the form
\begin{equation}
[\mathbf{K}_0,H_{\mathrm{int}}]=-[\mathbf{W},H].
\tag{3.3.21}\label{eq:3.3.21}
\end{equation}
By itself, the condition \eqref{eq:3.3.21} is empty, because for any $H_{\mathrm{int}}$ we
could always define $\mathbf{W}$ by giving its matrix elements between
$H$-eigenstates $\ket{\alpha}$ and $\ket{\beta}$ as
$-\bra{\beta}[\mathbf{K}_0,H_{\mathrm{int}}]\ket{\alpha}/
(E_\beta-E_\alpha)$. Recall that the crucial point in the Lorentz invariance of a
theory is not that there should exist a set of exact generators satisfying
Eqs.~\eqref{eq:3.3.11}--\eqref{eq:3.3.17}, but rather that these operators should act the same way on
`in' and `out' states; merely finding an operator $\mathbf{K}$ that satisfies
Eq.~\eqref{eq:3.3.21} is not enough. Eq.~\eqref{eq:3.3.21} does become significant if we add the
requirement that matrix elements of $\mathbf{W}$ should be smooth functions of the
energies, and in particular should not have singularities of the form
$(E_\beta-E_\alpha)^{-1}$. We shall now show that Eq.~\eqref{eq:3.3.21}, together with an
appropriate smoothness condition on $\mathbf{W}$, does imply the remaining Lorentz
invariance condition $[\mathbf{K}_0,S]=0$.

To prove this, let us consider the commutator of $\mathbf{K}_0$ with the operator
$U(t,t_0)$ for finite $t$ and $t_0$. Using Eq.~\eqref{eq:3.3.10} and the fact that
$\mathbf{P}_0$ commutes with $H_0$ yields:
\[
[\mathbf{K}_0,\exp(iH_0t)]=-t\mathbf{P}_0\exp(iH_0t)
\]
while Eq.~\eqref{eq:3.3.21} (which is equivalent to Eq.~\eqref{eq:3.3.17}) yields
\[
[\mathbf{K},\exp(iHt)]=-t\mathbf{P}\exp(iHt)
=-t\mathbf{P}_0\exp(iHt).
\]
The momentum operators then cancel in the commutator of $\mathbf{K}_0$ with $U$,
and we find:
\begin{equation}
[\mathbf{K}_0,U(\tau,\tau_0)]
=-\mathbf{W}(\tau)U(\tau,\tau_0)+U(\tau,\tau_0)\mathbf{W}(\tau_0),
\tag{3.3.22}\label{eq:3.3.22}
\end{equation}
where
\begin{equation}
\mathbf{W}(t)=\exp(iH_0t)\mathbf{W}\exp(-iH_0t).
\tag{3.3.23}\label{eq:3.3.23}
\end{equation}
If the matrix elements of $\mathbf{W}$ between $H_0$-eigenstates are sufficiently
smooth functions of energy, then matrix elements of $\mathbf{W}(t)$ between smooth
superpositions of energy eigenstates vanish for $t\to\pm\infty$, so Eq.~\eqref{eq:3.3.22}
gives in effect:
\begin{equation}
0=[\mathbf{K}_0,U(\infty,-\infty)]=[\mathbf{K}_0,S],
\tag{3.3.24}\label{eq:3.3.24}
\end{equation}
as was to be shown. This is the essential result: Eq.~\eqref{eq:3.3.21} together with the
smoothness condition on matrix elements of $\mathbf{W}$ that ensures that
$\mathbf{W}(t)$ vanishes for $t\to\pm\infty$ provides a sufficient condition for
the Lorentz invariance of the $S$-matrix. This smoothness condition is a natural
one, because it is much like the condition on matrix elements of $H_{\mathrm{int}}$
that is needed to make $H_I(t)$ vanish for $t\to\pm\infty$, as required in order
to justify the very idea of an $S$-matrix.

We can also use Eq.~\eqref{eq:3.3.22} with $\tau=0$ and $\tau_0=\mp\infty$ to show that
\begin{equation}
\mathbf{K}\Omega(\mp\infty)=\Omega(\mp\infty)\mathbf{K}_0,
\tag{3.3.25}\label{eq:3.3.25}
\end{equation}
where $\Omega(\mp\infty)$ is according to \eqref{eq:3.1.13} the operator that converts a
free-particle state $\ket{\alpha}_0$ into the corresponding `in' or `out' state.
Also, it follows trivially from Eqs.~\eqref{eq:3.3.18} and \eqref{eq:3.3.19} that the same is true
for the momentum and angular momentum:
\begin{align}
\mathbf{P}\Omega(\mp\infty)&=\Omega(\mp\infty)\mathbf{P}_0,
&\tag{3.3.26}\label{eq:3.3.26}\\
\mathbf{J}\Omega(\mp\infty)&=\Omega(\mp\infty)\mathbf{J}_0.
&\tag{3.3.27}\label{eq:3.3.27}
\end{align}
Finally, since all free and asymptotic states are eigenstates of $H_0$ and $H$
respectively with the same eigenvalue $E_\alpha$, we have
\begin{equation}
H\Omega(\mp\infty)=\Omega(\mp\infty)H_0.
\tag{3.3.28}\label{eq:3.3.28}
\end{equation}

Eqs.~\eqref{eq:3.3.25}--\eqref{eq:3.3.28} show that with our assumptions, `in' and `out' states do
transform under inhomogeneous Lorentz transformations just like the free-particle
states. Also, since these are similarity transformations, we now see that the exact
generators $\mathbf{K}$, $\mathbf{P}$, $\mathbf{J}$, and $H$ satisfy the same
commutation relations as $\mathbf{K}_0$, $\mathbf{P}_0$, $\mathbf{J}_0$, and $H_0$.
This is why it turned out to be unnecessary in proving the Lorentz invariance of
the $S$-matrix to use the other commutation relations \eqref{eq:3.3.12}, \eqref{eq:3.3.13}, and
\eqref{eq:3.3.15} that involve $\mathbf{K}$.

\subsubsection*{Internal Symmetries}

There are various symmetries, like the symmetry in nuclear physics under
interchange of neutrons and protons, or the `charge-conjugation' symmetry between
particles and antiparticles, that have nothing directly to do with Lorentz
invariance, and further appear the same in all inertial frames. Such a symmetry
transformation $T$ acts on the Hilbert space of physical states as a unitary
operator $U(T)$, that induces linear transformations on the indices labelling
particle species
\begin{align}
&U(T)\InOutKet{p_1,\sigma_1,n_1;p_2,\sigma_2,n_2;\ldots}
\nonumber\\
&\quad=\sum_{\bar n_1\bar n_2\ldots}
\mathcal{D}_{\bar n_1n_1}(T)\mathcal{D}_{\bar n_2n_2}(T)\cdots
\InOutKet{p_1,\sigma_1,\bar n_1;p_2,\sigma_2,\bar n_2;\ldots}.
\tag{3.3.29}\label{eq:3.3.29}
\end{align}
In accordance with the general discussion in Chapter 2, the $U(T)$ must satisfy
the group multiplication rule
\begin{equation}
U(\bar T)U(T)=U(\bar TT),
\tag{3.3.30}\label{eq:3.3.30}
\end{equation}
where $\bar TT$ is the transformation obtained by first performing the
transformation $T$, then some other transformation $\bar T$. Acting on
Eq.~\eqref{eq:3.3.29} with $U(\bar T)$, we see that the matrices $\mathcal{D}$ satisfy the
same rule
\begin{equation}
\mathcal{D}(\bar T)\mathcal{D}(T)=\mathcal{D}(\bar TT).
\tag{3.3.31}\label{eq:3.3.31}
\end{equation}
Also, taking the scalar product of the states obtained by acting with $U(T)$ on
two different `in' states or two different `out' states, and using the
normalization condition \eqref{eq:3.1.2}, we see that $\mathcal{D}(T)$ must be unitary
\begin{equation}
\mathcal{D}^\dagger(T)=\mathcal{D}^{-1}(T).
\tag{3.3.32}\label{eq:3.3.32}
\end{equation}
Finally, taking the scalar product of the states obtained by acting with $U(T)$
on one `out' state and one `in' state shows that $\mathcal{D}$ commutes with the
$S$-matrix, in the sense that
\begin{align}
&\sum_{\bar N_1\bar N_2\ldots}
\sum_{\bar N'_1\bar N'_2\ldots}
\mathcal{D}^*_{\bar N'_1n'_1}(T)\mathcal{D}^*_{\bar N'_2n'_2}(T)\cdots
\mathcal{D}_{\bar N_1n_1}(T)\mathcal{D}_{\bar N_2n_2}(T)\cdots
\nonumber\\
&\qquad\cdot
S_{p'_1\sigma'_1\bar N'_1;p'_2\sigma'_2\bar N'_2;\ldots,\,
p_1\sigma_1\bar N_1;p_2\sigma_2\bar N_2;\ldots}
\nonumber\\
&=S_{p'_1\sigma'_1n'_1;p'_2\sigma'_2n'_2;\ldots,\,
p_1\sigma_1n_1;p_2\sigma_2n_2;\ldots}.
\tag{3.3.33}\label{eq:3.3.33}
\end{align}

Again, this is a definition of what we mean by a theory being invariant under the
internal symmetry $T$, because to derive Eq.~\eqref{eq:3.3.33} we still need to show that
the \textit{same} unitary operator $U(T)$ will induce the transformation \eqref{eq:3.3.29}
on both `in' and `out' states. This will be the case if there is an `unperturbed'
transformation operator $U_0(T)$ that induces these transformations on
free-particle states,
\begin{align}
&U_0(T)\ket{p_1,\sigma_1,n_1;p_2,\sigma_2,n_2;\ldots}_0
\nonumber\\
&\quad=\sum_{\bar N_1\bar N_2\ldots}
\mathcal{D}_{\bar N_1n_1}(T)\mathcal{D}_{\bar N_2n_2}(T)\cdots
\ket{p_1,\sigma_1,\bar N_1;p_2,\sigma_2,\bar N_2;\ldots}_0
\tag{3.3.34}\label{eq:3.3.34}
\end{align}
and that commutes with both the free-particle and interaction parts of the
Hamiltonian
\begin{align}
U_0^{-1}(T)H_0U_0(T)&=H_0,
&\tag{3.3.35}\label{eq:3.3.35}\\
U_0^{-1}(T)H_{\mathrm{int}}U_0(T)&=H_{\mathrm{int}}.
&\tag{3.3.36}\label{eq:3.3.36}
\end{align}
From either the Lippmann--Schwinger equation \eqref{eq:3.1.17} or from \eqref{eq:3.1.13}, we see
that the operator $U_0(T)$ will induce the transformations \eqref{eq:3.3.29} on `in' and
`out' states as well as free-particle states, so that we can derive Eq.~\eqref{eq:3.3.29}
taking $U(T)$ as $U_0(T)$.

A special case of great physical importance is that of a one-parameter Lie group,
where $T$ is a function of a single parameter $\theta$, with
\begin{equation}
T(\bar\theta)T(\theta)=T(\bar\theta+\theta).
\tag{3.3.37}\label{eq:3.3.37}
\end{equation}
As shown in Section~2.2, in this case the corresponding Hilbert-space operators
must take the form
\begin{equation}
U(T(\theta))=\exp(iQ\theta)
\tag{3.3.38}\label{eq:3.3.38}
\end{equation}
with $Q$ a Hermitian operator. Likewise the matrices $\mathcal{D}(T)$ take the form
\begin{equation}
\mathcal{D}_{n'n}(T(\theta))=\delta_{n'n}\exp(iq_n\theta),
\tag{3.3.39}\label{eq:3.3.39}
\end{equation}
where $q_n$ are a set of real species-dependent numbers. Here Eq.~\eqref{eq:3.3.33} simply
tells us that the $q$s are conserved: $S_{\beta\alpha}$ vanishes unless
\begin{equation}
q_{n'_1}+q_{n'_2}+\cdots=q_{n_1}+q_{n_2}+\cdots.
\tag{3.3.40}\label{eq:3.3.40}
\end{equation}
The classic example of such a conservation law is that of conservation of electric
charge. Also, all known processes conserve baryon number (the number of baryons,
such as protons, neutrons, and hyperons, minus the number of their antiparticles)
and lepton number (the number of leptons, such as electrons, muons, $\tau$ particles,
and neutrinos, minus the number of their antiparticles) but as we shall see in
Volume II, these conservation laws are believed to be only very good approximations.
There are other conservation laws of this type that are definitely only approximate,
such as the conservation of the quantity known as strangeness, which was introduced
to explain the relatively long life of a class of particles discovered by Rochester
and Butler~\hyperref[ref:3.5]{[5]} in cosmic rays in 1947. For instance, the mesons now called $K^+$ and
$K^0$\footnote{Superscripts denote electric charges in units of the absolute value
of the electronic charge. A `hyperon' is any particle carrying non-zero strangeness
and unit baryon number.} are assigned strangeness $+1$ and the hyperons
$\Lambda^0$, $\Sigma^+$, $\Sigma^0$, $\Sigma^-$ are assigned strangeness $-1$,
while the more familiar protons, neutrons, and $\pi$ mesons (or pions) are taken to
have strangeness zero. The conservation of strangeness in strong interactions
explains why strange particles are always produced in association with one another,
as in reactions like $\pi^++n\to K^++\Lambda^0$, while the relatively slow decays
of strange particles into non-strange ones such as $\Lambda^0\to p+\pi^-$ and
$K^+\to\pi^++\pi^0$ show that the interactions that do not conserve strangeness
are very weak.

The classic example of a `non-Abelian' symmetry whose generators do not commute
with one another is isotopic spin symmetry, which was suggested~\hyperref[ref:3.6]{[6]} in 1937 on the
basis of an experiment~\hyperref[ref:3.7]{[7]} that showed the existence of a strong proton--proton
force similar to that between protons and neutrons. Mathematically, the group is
$SU(2)$, like the covering group of the group of three-dimensional rotations; its
generators are denoted $t_i$, with $i=1,2,3$, and satisfy a commutation relation
like \eqref{eq:2.4.18}:
\[
[t_i,t_j]=i\epsilon_{ijk}t_k.
\]
To the extent that isotopic spin symmetry is respected, it requires particles to
form degenerate multiplets labelled with an integer or half-integer $T$ and with
$2T+1$ components distinguished by their $t_3$ values, just like the degenerate
spin multiplets required by rotational invariance. These include the nucleons $p$
and $n$ with $T=\tfrac12$ and $t_3=\tfrac12,-\tfrac12$; the pions $\pi^+$,
$\pi^0$, and $\pi^-$ with $T=1$ and $t_3=+1,0,-1$; and the $\Lambda^0$ hyperon
with $T=0$ and $t_3=0$. These examples illustrate the relation between electric
charge $Q$, the third component of isotopic spin $t_3$, the baryon number $B$, and
the strangeness $S$:
\[
Q=t_3+(B+S)/2.
\]
This relation was originally inferred from observed selection rules, but it was
interpreted~\hyperref[ref:3.8]{[8]} by Gell-Mann and Ne'eman in 1960 to be a consequence of the
embedding of both the isospin $\mathbf{T}$ and the `hypercharge' $Y\equiv B+S$ in
the Lie algebra of a larger but more badly broken non-Abelian internal symmetry,
based on the non-Abelian group $SU(3)$. As we will see in Volume II, today both
isospin and $SU(3)$ symmetry are understood as incidental consequences of the small
masses of the two or three lightest quarks in the modern theory of strong
interactions, quantum chromodynamics.

The implications of isotopic spin symmetry for reactions among strongly interacting
particles can be worked out by the same familiar methods that were invented for
deriving the implications of rotational invariance. In particular, for a two-body
reaction $A+B\to C+D$, Eq.~\eqref{eq:3.3.33} requires that the $S$-matrix may be put in the
form (suppressing all but isospin labels):
\[
S_{t_{C3}t_{D3},t_{A3}t_{B3}}
=\sum_{T,t_3}C_{T_CT_D}(Tt_3;t_{C3}t_{D3})
C_{T_AT_B}(Tt_3;t_{A3}t_{B3})S_T,
\]
where $C_{j_1j_2}(j\sigma;\sigma_1\sigma_2)$ is the usual Clebsch--Gordan
coefficient~\hyperref[ref:3.9]{[9]} for forming a spin $j$ with three-component $\sigma$ from spins
$j_1$ and $j_2$ with three-components $\sigma_1$ and $\sigma_2$, respectively;
and $S_T$ is a `reduced' $S$-matrix depending on $T$ and on all the suppressed
momentum and spin variables, but not on the isospin three-components
$t_{A3},t_{B3},t_{C3},t_{D3}$. Of course, this, like all of the consequences of
isotopic spin invariance, is only approximate, because this symmetry is not
respected by electromagnetic (and other) interactions, as shown for instance by
the fact that different members of the same isospin multiplet like $p$ and $n$
have different electric charges and slightly different masses.
